\documentclass[aspectratio=169,10pt]{beamer}

\usetheme{default}
\setbeamertemplate{navigation symbols}{}
\setbeamersize{text margin left=0.65cm,text margin right=0.65cm}

\usepackage[T1]{fontenc}
\usepackage{lmodern}
\usepackage{graphicx}
\usepackage{booktabs}
\usepackage{array}
\usepackage{amsmath}
\usepackage{url}

\setbeamercolor{normal text}{fg=black,bg=white}
\setbeamercolor{frametitle}{fg=black,bg=white}
\setbeamerfont{frametitle}{series=\bfseries,size=\Large}
\setbeamertemplate{itemize item}{--}
\setbeamertemplate{itemize subitem}{--}
\setbeamertemplate{footline}{%
  \hfill{\scriptsize Compact magnetometer update \insertframenumber/\inserttotalframenumber}\hspace{0.6em}\vspace{0.35em}
}

\newcommand{\assetdir}{assets/project_update_20260511}
\newcommand{\smallitemsep}{\setlength{\itemsep}{0.45em}}

\title{Compact Magnetometer Project Update}
\subtitle{Lock-in module progress}
\author{Chalach Kasemtantikul}
\date{May 2026}

\begin{document}

\begin{frame}{Compact magnetometer project update}
\begin{columns}[T,totalwidth=\textwidth]
\begin{column}{0.56\textwidth}
\begin{itemize}\smallitemsep
  \item The lock-in module now runs the full workflow:
  \[
  \begin{aligned}
  \text{ODMR sweep} &\rightarrow \text{16 parked frequencies} \\
  &\rightarrow \Delta f(t) \rightarrow \Delta B(t).
  \end{aligned}
  \]
  \item Latest parking run: 8 kept transitions, 16 parked frequencies, no rejected fitted centers.
  \item Live lock-in test: 460 batches in 39.4 s, about 11.7 Hz update rate.
  \item Magnetic-field reconstruction is implemented with Yizhou's vector algorithm: 460/460 samples rank-ok.
  \item Live B-field measurement works, but online reconstruction slows the effective update rate.
\end{itemize}
\end{column}
\begin{column}{0.40\textwidth}
\centering
\begin{tabular}{@{}ll@{}}
\toprule
Item & Current result \\
\midrule
Transitions & 8 \\
Parked points & 16 \\
Single batch & 113 ms \\
Live lock-in & 11.7 Hz \\
Live B-field & 5.9 Hz \\
\bottomrule
\end{tabular}
\end{column}
\end{columns}
\end{frame}

\begin{frame}{Workflow 1: choose 16 parked frequencies}
\begin{columns}[T,totalwidth=\textwidth]
\begin{column}{0.70\textwidth}
\centering
\includegraphics[width=\linewidth,height=0.72\textheight,keepaspectratio]{\assetdir/Lockin_module_cell011_out04.png}
\end{column}
\begin{column}{0.27\textwidth}
\begin{itemize}\smallitemsep
  \item Start from the newest ODMR sweep.
  \item Fit 8 transition centers.
  \item Put two parked frequencies on the flanks of each transition.
  \item These flank points convert lock-in intensity changes into a frequency shift.
\end{itemize}
\end{column}
\end{columns}
\end{frame}

\begin{frame}{Workflow 2: acquire all parked points}
\begin{columns}[T,totalwidth=\textwidth]
\begin{column}{0.72\textwidth}
\centering
\includegraphics[width=\linewidth,height=0.72\textheight,keepaspectratio]{\assetdir/Lockin_module_cell013_out05.png}
\end{column}
\begin{column}{0.25\textwidth}
\begin{itemize}\smallitemsep
  \item One FPGA program measures all 16 frequencies.
  \item The output is one wide row per batch: signal, reference, and frequency for every parked point.
  \item Latest one-batch acquisition took 113 ms.
\end{itemize}
\end{column}
\end{columns}
\end{frame}

\begin{frame}{Workflow 3: convert lock-in signal to frequency shift}
\begin{columns}[T,totalwidth=\textwidth]
\begin{column}{0.49\textwidth}
\footnotesize
For each transition, use the two flanks:
\[
\Delta f_k(t)=
\frac{\left[S_+(t)-S_-(t)\right]-\left[S_{+,0}-S_{-,0}\right]}
{m_- - m_+}.
\]
\vspace{0.2em}
\begin{itemize}\setlength{\itemsep}{0.25em}
  \item $S_-$ and $S_+$ are the live signals at the lower and upper parked points.
  \item $S_{-,0}$ and $S_{+,0}$ come from the reference ODMR calibration.
  \item $m_-$ and $m_+$ are the local ODMR slopes at those two parked points.
  \item This gives a frequency-shift time series for each of the 8 transitions.
  \item Latest live run: 460 batches over 39.4 s.
  \item Per-block frequency noise: 0.08 to 0.25 MHz standard deviation.
\end{itemize}
\end{column}
\begin{column}{0.48\textwidth}
\centering
\includegraphics[width=\linewidth,height=0.70\textheight,keepaspectratio]{\assetdir/Lockin_module_cell021_out03.png}
\end{column}
\end{columns}
\end{frame}

\begin{frame}{Workflow 4: inferred new peak frequencies}
\begin{columns}[T,totalwidth=\textwidth]
\begin{column}{0.66\textwidth}
\centering
\includegraphics[width=\linewidth,height=0.78\textheight,keepaspectratio]{\assetdir/Lockin_module_cell021_out04.png}
\end{column}
\begin{column}{0.31\textwidth}
\begin{itemize}\smallitemsep
  \item This is the inferred new peak frequency for each of the 8 transitions.
  \item Solid traces show $f_{\mathrm{new}}(t)=f_0+\Delta f(t)$.
  \item Dashed lines mark the reference peak positions from the ODMR calibration.
  \item This view is useful for spotting sudden jumps before B-field reconstruction.
\end{itemize}
\end{column}
\end{columns}
\end{frame}

\begin{frame}{Workflow 5: reconstruct magnetic field}
\begin{columns}[T,totalwidth=\textwidth]
\begin{column}{0.76\textwidth}
\centering
\includegraphics[width=\linewidth,height=0.49\textheight,keepaspectratio]{\assetdir/Lockin_module_cell022_out04.png}

\vspace{0.15em}
\includegraphics[width=\linewidth,height=0.23\textheight,keepaspectratio]{\assetdir/Lockin_module_cell022_out05.png}
\end{column}
\begin{column}{0.21\textwidth}
\scriptsize
\begin{itemize}\setlength{\itemsep}{0.18em}
  \item 8 projected shifts to vector B-field.
  \item Offline live-run processing: 460 samples, 460 rank-ok.
  \item Projected-field noise: 2.9 to 8.9 $\mu$T.
  \item Reconstructed B statistics:
\end{itemize}
\vspace{0.1em}
\centering
\tiny
\begin{tabular}{@{}lrr@{}}
\toprule
Comp. & Mean & Std. \\
\midrule
$\Delta B_x$ & -3.92 & 3.54 \\
$\Delta B_y$ & -13.89 & 4.69 \\
$\Delta B_z$ & -3.07 & 4.09 \\
$|\Delta B|$ & 15.79 & 4.43 \\
\bottomrule
\end{tabular}

\vspace{0.15em}
\tiny Units: $\mu$T
\end{column}
\end{columns}
\end{frame}

\begin{frame}{Workflow 6: live B-field measurement}
\begin{columns}[T,totalwidth=\textwidth]
\begin{column}{0.74\textwidth}
\centering
\includegraphics[width=\linewidth,height=0.82\textheight,keepaspectratio]{\assetdir/Lockin_module_cell024_out07.png}
\end{column}
\begin{column}{0.23\textwidth}
\scriptsize
\begin{itemize}\setlength{\itemsep}{0.15em}
  \item Same reconstruction path as the offline workflow.
  \item Latest live B run: 208 batches in 35.4 s.
  \item Update rate: about 5.9 Hz.
  \item Main bottleneck: per-batch reconstruction and plotting.
\end{itemize}
\vspace{0.2em}
\centering
\tiny
\begin{tabular}{@{}lrr@{}}
\toprule
Comp. & Mean & Std. \\
\midrule
$\Delta B_x$ & -4.26 & 3.91 \\
$\Delta B_y$ & -13.81 & 4.40 \\
$\Delta B_z$ & -3.05 & 4.98 \\
$|\Delta B|$ & 16.22 & 3.80 \\
\bottomrule
\end{tabular}

\vspace{0.15em}
\scriptsize Units: $\mu$T
\end{column}
\end{columns}
\end{frame}

\begin{frame}{Next 1: stabilize live traces and reject jumps}
\begin{columns}[T,totalwidth=\textwidth]
\begin{column}{0.32\textwidth}
\centering
\includegraphics[height=0.80\textheight,keepaspectratio]{\assetdir/Initial_Start_Live_Test_cell015_out02.png}
\end{column}
\begin{column}{0.32\textwidth}
\centering
\includegraphics[height=0.80\textheight,keepaspectratio]{\assetdir/Noisy_Live_Test_After_Booting_cell015_out02.png}
\end{column}
\begin{column}{0.32\textwidth}
\footnotesize
\begin{itemize}\setlength{\itemsep}{0.3em}
  \item Initial run: channels were relatively stable, but slowly decayed downward.
  \item Later runs: some parked channels developed sudden jumps.
  \item Short-term fix: reject sudden outliers before B reconstruction.
  \item Algorithm fix: add stricter checks so parked points avoid bad ODMR regions.
  \item Hardware diagnosis: separate warm-up/baseline drift from real peak jumps.
\end{itemize}
\end{column}
\end{columns}
\end{frame}

\begin{frame}{Next 2: verify reconstruction and speed up live B}
\begin{columns}[T,totalwidth=\textwidth]
\begin{column}{0.48\textwidth}
\textbf{Verify B reconstruction}
\begin{itemize}\smallitemsep
  \item Apply a known test field and compare the recovered direction and magnitude.
  \item Compare offline reconstruction and live B reconstruction on the same saved lock-in data.
  \item Track residual norm, matrix rank, and condition number for every batch.
\end{itemize}
\end{column}
\begin{column}{0.48\textwidth}
\textbf{Make live B faster}
\begin{itemize}\smallitemsep
  \item Decouple acquisition from reconstruction so data collection does not wait on plotting.
  \item Vectorize the per-batch calculation where possible.
  \item Update plots more slowly than the acquisition loop.
  \item Target: move live B closer to the raw live lock-in rate.
\end{itemize}
\end{column}
\end{columns}
\end{frame}

\begin{frame}{Next 3: system integration}
\begin{itemize}\smallitemsep
  \item Dual-channel measurement: connect the lock-in workflow to the dual-channel PL readout.
  \item Hardware optimization: stabilize warm-up behavior, ADC baseline, microwave timing, and optical alignment.
  \item Data organization: keep the GitHub workflow clean so Yizhou's reconstruction code and the live data stay synchronized.
  \item Patent manuscript: document the 16-frequency parked lock-in workflow, live peak tracking, and vector B-field reconstruction.
\end{itemize}
\end{frame}

\end{document}
